\section{Introduction}

Modern regular-expression engines are not merely membership testers for regular
languages.  A JavaScript engine must return the match selected by an ordered
backtracking policy, including the effects of alternation priority, greedy and
lazy quantifiers, captures, lookaround, and other language features.  Recent
work by Barriere, Deng, and Pit-Claudel gives a mechanized semantics for
JavaScript regular expressions in which the result of matching is not just a
single top-priority match, but a complete ordered backtracking tree recording the
matches and failures considered by the engine~\cite{barriere2026jsregex}.  That
tree is an attractive semantic object for verification: it is faithful to
backtracking while still being inductive and suitable for proof.

At the same time, the classical theory of finite automata has long studied the
number of accepting runs of a nondeterministic finite automaton (NFA) on a word.
This number is the word's ambiguity, and the supremum over all words is the
degree of ambiguity of the automaton.  Weber and Seidl showed that ambiguity is
not just a coarse counting measure: finite ambiguity admits structural bounds,
infinite ambiguity has a simple polynomial-time criterion, and the polynomial
degree of ambiguity growth is computable~\cite{weber1991ambiguity}.  These
results are tantalizing for ReDoS, because catastrophic backtracking is precisely
a phenomenon in which a compact regular expression induces a very large search
space on a family of inputs.

This paper asks whether these two accounts can be made to meet:
\begin{quote}
Can the growth of JavaScript-style backtracking traces be explained, bounded, and
eventually certified through automata ambiguity?
\end{quote}

Our answer is a formal mechanism that connects successful leaves of a
backtracking trace to accepting runs of a position automaton.  The intended
scope is deliberately staged.  We first treat the regular core of JavaScript
regular expressions, where Glushkov position automata give a syntax-directed NFA
whose states are occurrences of atoms.  We then use the tree semantics of
Barriere et al. as the target view of engine execution.  The central theorem
states that, after erasing priority annotations and failing leaves, successful
trace leaves are in bijection with accepting marked runs of the position
automaton.  This theorem is the bridge that allows ambiguity results to be
transported to backtracking traces.

The bridge gives three payoffs.

\begin{itemize}
\item It separates \emph{semantic nondeterminism} from \emph{engine priority}.
Ambiguity counts the possible successful decompositions of a word; the
backtracking policy orders those decompositions and selects the first.
\item It explains why unambiguity and finite ambiguity are useful ReDoS
invariants.  If the constructed automaton is unambiguous or finitely ambiguous,
then the successful part of the trace tree is correspondingly bounded.
\item It identifies the next verification target: lift Weber--Seidl criteria
into Coq as certificates for infinite or polynomial ambiguity growth, and use
those certificates to justify static ReDoS warnings.
\end{itemize}

This is an early-stage ICFP draft.  The Coq development already contains a
substantial regular-expression and automata core, but the full bridge theorem
and the Weber--Seidl classification proof are not complete.  We state them as
the main technical obligations of the paper rather than as finished results.

\paragraph{Contributions.}
We make the following contributions.

\begin{itemize}

\item
\textbf{A string-constraint semantics for positive lookahead.}
We introduce a language model for regular expressions with positive
lookahead in which a match is represented by a pair
$(u,v)\in\Sigma^*\times\Sigma^*$. The first component $u$ records the
string actually consumed, while the second component $v$ records the
constraint imposed on the following input. This separates real
consumption from zero-width lookahead requirements and provides the
semantic basis for reasoning about lookahead-sensitive matching.

\item
\textbf{An algebra of string-constraint languages.}
We define a prefix-compatible concatenation operation $\odot$ for
languages over $\Sigma^*\times\Sigma^*$. The operation propagates
right-context constraints through sequential composition by combining
residual consumption with context joining. We prove associativity,
identity laws, distributivity over union, and zero laws, showing that
string-constraint languages form an idempotent semiring. We then define
the Kleene star induced by $\odot$ and characterize the rational
string-constraint languages generated by expressions.

\item
\textbf{A partial-derivative continuation automaton.}
We develop left quotients, derivatives, and partial derivatives for
string-constraint languages, and refine partial derivatives by marked
consuming occurrences into $c$-derivatives and continuations. Based on
this residual calculus, we construct a lookahead-aware continuation
automaton whose states are derivative continuations rather than
position states. Transitions are induced by $c$-derivatives, while
lookahead assertions are treated as zero-width constraints or witness
conditions.

\item
\textbf{A Rocq mechanization of the semantics and automaton correctness.}
We formalize the string-constraint semantics, prefix-context
concatenation, residual operations, $c$-derivatives, continuations, and
the continuation automaton in Rocq. We prove that the automaton is sound
and complete with respect to the source expression semantics: the
language recognized by the automaton coincides with the
string-constraint language denoted by the expression. This gives a
machine-checked correctness theorem for positive-lookahead regular
expressions under the proposed derivative-continuation construction.

\end{itemize}


